\documentclass[12pt]{article}

\usepackage[margin=1in]{geometry}
\usepackage[T1]{fontenc}
\usepackage[utf8]{inputenc}
\usepackage{lmodern}
\usepackage{amsmath,amssymb,amsthm}
\usepackage{booktabs}
\usepackage{array}
\usepackage{float}
\usepackage{enumitem}
\usepackage{setspace}
\usepackage{natbib}
\usepackage{microtype}
\usepackage[colorlinks=true,linkcolor=blue,citecolor=blue,urlcolor=blue]{hyperref}

\onehalfspacing
\setlist[itemize]{topsep=3pt,itemsep=2pt,parsep=0pt,leftmargin=18pt}
\setlist[enumerate]{topsep=3pt,itemsep=2pt,parsep=0pt,leftmargin=18pt}

\newtheorem{proposition}{Proposition}
\newtheorem{assumption}{Assumption}
\newcolumntype{L}[1]{>{\raggedright\arraybackslash}p{#1}}
\newcommand{\E}{\mathbb{E}}
\newcommand{\V}{\mathbb{V}}
\newcommand{\ind}{\mathbb{1}}
\newcommand{\Rhat}{\widehat{R}}
\newcommand{\That}{\widehat{T}}
\newcommand{\Yba}{\bar{Y}}
\newcommand{\Npop}{N}
\newcommand{\sample}{\mathcal{S}}
\newcommand{\universe}{\mathcal{U}}

\title{Sampling Streets:\\
Design-Based Inference for Street-Level Field Measurement}
\author{Varun K.R. \and Gaurav Sood\thanks{Replication materials, including the spatial sampler, the routing allocator, the GoPro and Street View imagery pipelines, the sound meter pipeline, and the inference package, are at \url{https://github.com/geosensing}.}}
\date{May 2026}

\begin{document}
\maketitle

\begin{abstract}

Researchers increasingly use street-level field measurements  to study policy and political outcomes, from infrastructure quality and ambient and sound pollution to gendered public presence. Often, the data come from a convenience corpus rather than a probability sample: access determines where data are collected, and access often correlates with the outcomes measured. As a result, estimates are reported without a clearly defined estimand or a design-based standard error.

In this paper, we develop and introduce a probability-sampling framework for street-level field measurement.
This framework constructs a sampling frame from the road network, allocates sampled units into field itineraries or independent random walks, and produces design-based estimates with variance estimators for simple random, stratified, clustered, and walk-based designs. Acquisition modes include researcher-collected imagery, researcher-deployed sensors such as sound meters, and platform imagery such as Google Street View; each is treated as a measurement instrument with its own coverage, time, and bias properties. Platform imagery faces three measurement-relevant constraints in particular: obfuscation of faces and plates, collection at platform-determined times rather than researcher-randomized times, and image dates that may lag the question by years. We illustrate with two applications: women's presence in public space across Mumbai, Navi Mumbai, and Delhi using both researcher-collected and platform imagery, and street-level noise exposure across Delhi using sound meters. The framework brings image- and sensor-based political measurement into the standard machinery of survey inference.
\end{abstract}

\section{Introduction}

Field measurements at sampled street locations are an increasingly common source of data for political and social science. Researchers measure neighborhood disorder, sidewalk infrastructure, traffic conditions, urban inequality, and visible composition of public space from street-level imagery \citep{rundle2011streetview,kelly2013streetview,naik2014streetscore,gebru2017streetview}, and they measure environmental quantities such as noise, particulate matter, and temperature from spatially distributed sensors. Most of this work draws on data that are convenient: large platform corpora such as Google Street View, single-city censuses of named blocks, sensor placements at fixed monitoring stations, or routes the researcher can reach. The convenience corpus delivers volume, but it does not deliver an estimand. The resulting numbers describe what is in the corpus rather than what is in the city, and standard errors, when reported, typically treat units as independent draws from a hypothetical superpopulation rather than as design-based samples from the actual road network.

We develop a probability-sampling framework for street-level field measurement. The frame is constructed from the road network. Sampling is randomized over street segments or independent random walks. Allocation to field collection is treated as an operational step, not a sampling step. Estimation is design-based. The framework supports three acquisition modes: researcher-collected imagery, suitable when visual outcomes vary with time of day and current conditions; researcher-deployed sensors, suitable when the quantity of interest is a continuous physical measurement; and platform imagery, suitable when outcomes change slowly and the platform's coverage and date metadata are appropriate to the question. Coverage and nonresponse are part of the design rather than a robustness exercise.

We illustrate the framework with two applications. The first is women's presence in public space, an outcome motivated by the long-running concern that women are underrepresented in everyday public life \citep{sen1990missing,phadke2011whyloiter}. The second is street-level noise exposure in Delhi, an outcome motivated by environmental-health evidence that road-traffic noise affects sleep, cardiovascular health, cognition, and annoyance \citep{whoGuidelines2018,whoCompendium2022}. The two applications use different acquisition modes (imagery and sound meters), produce different estimands (a population ratio and a population mean), and share the same sampling and inference machinery.

The methodological contribution is to bring image- and sensor-based political measurement into the standard machinery of survey inference. The applied contribution is a pair of design-based estimates that are explicitly tied to the road network sampled, the time windows during which collection occurred, and the measurement instrument used.

\section{Estimands}

\subsection{General setup}

Let $i \in \universe = \{1, \ldots, \Npop\}$ index finite population units. A unit may be a road segment, a camera heading at a road segment, a sound meter reading at a road-segment-by-time cell, or a sampled point along an allocated route. For unit $i$, let $Y_i$ denote a numeric measurement and, when relevant, $Z_i$ a corresponding denominator measurement. The interpretation of $Y_i$ and $Z_i$ depends on the application. For street imagery counting visible people by gender, $Y_i = W_i$ is the number of women visible and $Z_i = H_i = W_i + M_i$ is the number of classified persons. For sound measurement, $Y_i$ is a continuous reading (e.g., the maximum decibel level within a fixed measurement window) and no denominator is used.

\subsection{Three estimands}

Three estimands cover the applications considered here.

The \emph{population ratio} is
\begin{equation}
    R = \frac{\sum_{i \in \universe} Y_i}{\sum_{i \in \universe} Z_i}.
\end{equation}
For the imagery application, $R$ is the people-weighted female share: the share of all visible people in the target population who are annotated as women. It weights each location by the number of people present, so a market with two hundred visible people contributes more than a quiet side street with two.

The \emph{population mean} is
\begin{equation}
    \mu = \frac{1}{\Npop} \sum_{i \in \universe} Y_i.
\end{equation}
For the noise application, $\mu$ is the expected sound level at a randomly drawn unit in the target population. It is the natural design-based analogue of ``how loud is the average street'' when the unit is a sampled point on the road network.

The \emph{population threshold-exceedance share} is
\begin{equation}
    P_\tau = \frac{1}{\Npop} \sum_{i \in \universe} \ind\{Y_i \geq \tau\}.
\end{equation}
$P_\tau$ answers ``what share of sampled units register at or above $\tau$,'' a useful estimand for outcomes with regulatory thresholds. For the noise application, $\tau = 55$ dB corresponds to the daytime residential ambient standard under India's Noise Pollution Rules, and $P_{55}$ has a direct policy interpretation.

A fourth quantity, the location-weighted mean of a per-unit ratio,
\begin{equation}
    L = \frac{1}{\Npop^+} \sum_{i \in \universe: Z_i > 0} \frac{Y_i}{Z_i},
\end{equation}
where $\Npop^+$ counts nonempty units, is the mean of a per-unit female share across nonempty street locations. $L$ weights a quiet side street with two people the same as a market with two hundred. Reporting $R$ and $L$ together is informative because their divergence carries content. A higher $R$ than $L$ means that busier locations have proportionally more women; a lower $R$ than $L$ means that high-traffic locations are especially male dominated.

A fifth quantity ties any of the above to a benchmark. For the imagery application, let $B$ be a chosen baseline female share derived from census sex ratios, age-adjusted local populations, labor-force participation, transit ridership, or a matched-city comparison. The public-presence gap is $G = R - B$; a negative $G$ means that women are underrepresented in public space relative to the benchmark.

\section{Sampling Design}

\subsection{The spatial frame}

Let the city road network be represented as a finite collection of eligible street segments. The eligible network is defined by road classes, such as primary, secondary, tertiary, residential, and unclassified roads under the OpenStreetMap classification \citep{openstreetmap}. Roads are split into approximately half-kilometer segments, segment coordinates are recorded, and segment rows constitute the sampling frame. The frame is finite, enumerable, and reproducible from the OSM extract.

The frame defines what the estimate is about. A frame restricted to arterial roads answers a question about arterial roads; a frame that includes residential streets answers a question about the eligible road network as a whole. A stratified design across road types answers both, provided that design weights are used in estimation.

\subsection{Selection mechanisms}

Two selection mechanisms are useful. The first is point-based sampling: draw $n$ segments without replacement from the frame, with inclusion probabilities $\pi_i$. Equal-probability simple random sampling sets $\pi_i = n/\Npop$. Stratified sampling sets $\pi_{ih} = n_h/N_h$ within stratum $h$, where $N_h$ is the number of segments in stratum $h$ and $n_h$ is the number sampled.

The second mechanism is walk-based sampling: draw a sample of starting points and a length, generate independent random walks on the road network, and take measurements along each walk at a fixed cadence. The independent walks become the primary sampling units, and within-walk observations are systematic sub-samples. The walk design is useful when the substantive object is the experience of moving through the city rather than a static cross-section of segments, and it is the natural design for mobile sensors that record continuously along a route.

\subsection{Allocation as an operational layer}

Probability sampling chooses the locations; field logistics determine how those locations are visited. Sampled points are clustered into trip groups, ordered through a routing engine, and assigned to fieldworkers; for the walk design, walks are generated and assigned. Clustering and routing are operational steps that follow sampling. They affect the variance of the estimator, because adjacent points on the same itinerary may share neighborhoods, time windows, fieldworkers, and acquisition conditions, but they do not change the inclusion probabilities. Inference therefore proceeds from the original $\pi_i$, with itinerary structure used to construct cluster-robust standard errors and walk structure used to construct walk-based standard errors.

\section{Estimation and Inference}

\subsection{Horvitz-Thompson totals}

Let $S_i \in \{0, 1\}$ denote whether population unit $i$ is included in the sample. The Horvitz-Thompson estimators for population totals are
\begin{equation}
    \That_Y = \sum_{i \in \sample} \frac{Y_i}{\pi_i}, \qquad
    \That_Z = \sum_{i \in \sample} \frac{Z_i}{\pi_i}.
\end{equation}
Under the sampling design, $\That_Y$ and $\That_Z$ are unbiased for $T_Y = \sum_i Y_i$ and $T_Z = \sum_i Z_i$ \citep{horvitz1952generalization,lohr2019sampling}.

\subsection{Ratio estimator}

The ratio estimator for $R$ is
\begin{equation}
    \Rhat = \frac{\That_Y}{\That_Z}.
\end{equation}

\begin{proposition}[Design consistency and finite-sample bias]
\label{prop:consistency}
If the sampling design has known positive inclusion probabilities, the measurement process returns fixed values $(Y_i, Z_i)$ for sampled units, and standard regularity conditions for ratio estimation hold, then $\Rhat$ is design-consistent for $R$ and its bias is of order $O(n^{-1})$. Under simple random sampling without replacement, the leading variance is consistently estimated by a Taylor linearization using residuals $e_i = Y_i - \Rhat Z_i$.
\end{proposition}

The total estimators are unbiased; the ratio estimator is consistent with $O(n^{-1})$ finite-sample bias because it divides two random totals \citep{cochran1977sampling}. The bias term is small in samples typical of street-level studies and can be estimated directly from the data.

\subsection{Mean estimator}

When the estimand is a population mean rather than a ratio, the same machinery applies with $Z_i \equiv 1$. Under simple random sampling, the design-based estimator is the sample mean $\Yba = n^{-1} \sum_{i \in \sample} Y_i$, which is design-unbiased for $\mu$, and the variance estimator is
\begin{equation}
    \widehat{\V}(\Yba) \;=\; \left(1 - \frac{n}{\Npop}\right) \frac{1}{n} \cdot \frac{1}{n-1} \sum_{i \in \sample} (Y_i - \Yba)^2,
\end{equation}
with the finite population correction dropped when the sampling fraction is small. For the threshold-exceedance estimand, $Y_i$ is replaced by $\ind\{Y_i \geq \tau\}$ and the estimator reduces to the sample proportion above the threshold, with the corresponding Bernoulli variance.

\subsection{Variance under simple random sampling for the ratio}

For equal-probability simple random sampling of $n$ units from $\Npop$, $\pi_i = n/\Npop$, the weights cancel in the ratio, and
\begin{equation}
    \Rhat = \frac{\sum_{i \in \sample} Y_i}{\sum_{i \in \sample} Z_i}.
\end{equation}
Let $\bar{Z} = n^{-1} \sum_{i \in \sample} Z_i$ and $e_i = Y_i - \Rhat Z_i$. The linearized variance estimator is
\begin{equation}
    \widehat{\V}(\Rhat) \;=\; \left(1 - \frac{n}{\Npop}\right) \frac{1}{n \bar{Z}^2} \cdot \frac{1}{n-1} \sum_{i \in \sample} (e_i - \bar{e})^2,
\end{equation}
\citep{cochran1977sampling,lohr2019sampling}.

\subsection{Stratification}

When the eligible road network is partitioned into strata $h = 1, \ldots, H$ with $N_h$ segments and $n_h$ sampled, the design weight for sampled unit $i$ in stratum $h$ is $a_i = N_h / n_h$. The weighted ratio estimator is
\begin{equation}
    \Rhat_{st} = \frac{\sum_{i \in \sample} a_i Y_i}{\sum_{i \in \sample} a_i Z_i},
\end{equation}
and the corresponding stratified mean is $\Yba_{st} = (\sum a_i)^{-1} \sum a_i Y_i$. Stratification matters when road types differ in the outcome. The women's-presence application finds that the female share is lower on primary and secondary roads and higher in residential areas; the noise application finds that decibel levels vary with road class and traffic density. In both, an unweighted estimator from a road-type-imbalanced sample would be biased.

\subsection{Cluster-robust variance under itinerary allocation}

Allocator-created itineraries can introduce dependence among observations even when sampling is independent. Adjacent sampled points share neighborhoods, time windows, fieldworkers, weather, and acquisition conditions. A cluster-robust linearization aggregates residuals within each itinerary $g$:
\begin{equation}
    E_g = \sum_{i \in g \cap \sample} a_i (Y_i - \Rhat Z_i), \qquad
    \widehat{\V}_{cl}(\Rhat) = \frac{G}{G-1} \cdot \frac{\sum_{g=1}^{G}(E_g - \bar{E})^2}{\Big(\sum_{i \in \sample} a_i Z_i\Big)^2},
\end{equation}
where $G$ is the number of itineraries \citep{cameron2015cluster}. The corresponding cluster-robust variance for the mean replaces the residual with $a_i Y_i$ in the numerator and $\sum a_i$ in the denominator. The ratio of cluster-robust to naive standard errors is itself a useful diagnostic: a ratio near one indicates that operational clustering did not induce dependence in the outcome, while a large ratio indicates that itinerary structure carries substantial information.

\subsection{Walk-based variance}

When the design generates independent random walks, the walk is the primary sampling unit and inference is between walks. Let $Y_g = \sum_{i \in g} Y_i$ and $Z_g = \sum_{i \in g} Z_i$ aggregate counts within walk $g$. The ratio is $\Rhat = \sum_g Y_g / \sum_g Z_g$, and inference uses the variation across the $G$ independent walks. For a mean, $\Yba = \sum_g Y_g / \sum_g n_g$ where $n_g$ is the number of observations on walk $g$, and the variance is the corresponding between-walk variance. Under the walk design, the walk-based standard error is the design-preferred standard error, not a robustness check.

\subsection{Coverage and nonresponse}

A spatial sample is not complete until each sampled unit is either observed or coded as nonresponse. Let $C_i \in \{0, 1\}$ indicate whether usable data are available for unit $i$. For Google Street View, $C_i$ depends on whether platform vehicles have covered the area, whether the imagery is user-contributed, and whether the panorama is too old, blurred, or otherwise unusable. For field-collected imagery and sensor data, $C_i$ depends on worker compliance, road access, equipment failure, weather, and safety constraints.

If $C_i$ is independent of $(Y_i, Z_i)$ conditional on design strata $X_i$, complete-case weighted estimation is design-unbiased, and the inverse coverage probability multiplies the existing weight: $a_i^* = a_i / \widehat{\Pr}(C_i = 1 \mid X_i)$. If $C_i$ is associated with $(Y_i, Z_i)$ even after conditioning, the estimand changes from a population quantity to a quantity conditional on coverage, and the gap between coverage-adjusted and unadjusted estimates is itself part of the result.

\section{Acquisition Modes}

The framework supports three acquisition modes. The choice among them is governed by what the estimand requires from the time dimension, from the spatial coverage of available data, and from the physical quantity being measured.

\subsection{Researcher-collected imagery}

A field team carrying GoPro or equivalent cameras visits assigned itineraries, records video tagged with GPS and EXIF metadata, and uploads frames for annotation. A frame-to-GPS index assigns each annotated frame to a sampled coordinate, and the resulting counts feed the estimator. The mode preserves randomization over time of day, day of week, and weather, because the field team visits on the assigned schedule. The cost is operational: field collection is more expensive per unit than a platform query, and itinerary compliance has to be verified against GPS tracks rather than assumed.

The annotation produces counts of visible people by coarse categories. The categories are constructed from visible cues, not from observed identity. The protocol records ``unknown'' and ``ambiguous'' as explicit categories, and inference reports both a complete-case female share among classified persons and bounds in which unknowns are alternately classified as all women or all men. The bounds are tight when classification rates are high and informative when they are not.

\subsection{Researcher-deployed sensors}

A field team carries a calibrated sensor, follows assigned itineraries, and records measurements at a fixed cadence along the route and at substantively important conditions (such as traffic stops for noise measurement). Each observation stores GPS, timestamp, and the relevant physical reading. The sensor mode applies to any quantity that can be measured in situ at a sampled point: noise, particulate matter, temperature, illuminance, or other environmental signals.

The mode preserves the same randomization properties as researcher-collected imagery. Itinerary compliance is verified against GPS; observation timing is verified against the route schedule. The substantive estimand is typically a population mean or a threshold-exceedance share rather than a per-unit ratio, but the sampling and inference machinery is the same.

\subsection{Platform imagery}

Google Street View provides a metadata endpoint that returns panorama availability, location, panorama ID, image date, and copyright information \citep{googleStreetViewMetadata}. A Static API returns non-interactive panoramas at requested headings and requires an API key and billing \citep{googleStreetViewStatic}. The sampler queries metadata at sampled coordinates, downloads images where coverage exists, and feeds them to annotation. Platform imagery is scalable, low marginal cost, and well-suited to outcomes that change slowly.

Three measurement-relevant constraints distinguish platform imagery from researcher-collected modes:

\begin{enumerate}
\item \emph{Obfuscation.} Google applies automatic privacy protections that blur identifiable faces and license plates \citep{googleBlurPolicy,frome2009privacy}. Blurring is ethically valuable, but for any annotation that relies on facial cues (gender, age, expression), it raises the share of ambiguous cases and depresses the classified denominator. The measurement implication is that platform imagery is better suited to outcomes that do not require facial inference (presence vs absence of persons, infrastructure features, vehicles, vendors) than to outcomes that do.
\item \emph{Time of collection.} Capture time is set by the platform, not the researcher. Imagery is collected during the hours when platform vehicles operate and during the seasons in which they cover a given region. The resulting frame is biased toward daylight on dry weekdays and away from evenings, weekends, and weather conditions that are themselves substantively interesting. For outcomes that vary sharply by hour, day of week, or weather (public presence, market activity, traffic congestion), platform imagery does not estimate the time-randomized quantity that the researcher would obtain by visiting at randomized times.
\item \emph{Recency.} Image dates can be years out of date, and the granularity of the date metadata can vary or be missing \citep{curtis2013streetviewdates}. A panorama captured three years before the question of interest documents a different state of the city. For built-environment features that change slowly (footpaths, lane markings), this is often acceptable; for features that change quickly (signage, vendor activity, infrastructure investments), it is not.
\end{enumerate}

The three constraints do not disqualify platform imagery. They locate it. The mode is most appropriate for slowly-changing infrastructure outcomes, for spatially expansive audits where the marginal cost of researcher collection is prohibitive, and for studies in which the platform's observed time and date metadata are part of the analysis rather than ignored. For outcomes about current street-times and human presence, the framework treats platform imagery as a complement to researcher-collected modes rather than a substitute.

\section{Applications}

\subsection{Women's presence in public space}

The first application targets $R$, the people-weighted female share, in three Indian cities. Field collection in Mumbai and Navi Mumbai using the GoPro mode estimates women at roughly 18 to 19 percent of visible pedestrians, well below census sex-ratio benchmarks. Google Street View sampling across Mumbai, Delhi, and Navi Mumbai estimates women at roughly 21 to 24 percent of visible pedestrians, with the female share lower on primary and secondary roads and higher in residential areas. The two modes give similar overall magnitudes in the overlapping cities. The remaining gap is consistent with the platform-imagery time bias (Street View is captured during daytime hours when Google vehicles operate) and with the road-type composition of the underlying samples.

The estimates are reported with design-based confidence intervals around $R$, $L$, and $G$, with cluster-robust standard errors when itinerary allocation was used and walk-based standard errors when independent walks were the sampling unit. The two estimands differ for these data: $R$ is meaningfully lower than $L$, indicating that high-traffic locations are especially male dominated. The gap $G = R - B$ is large and negative against any plausible benchmark.

\subsection{Street-level noise exposure in Delhi}

The second application targets $\mu$ and $P_\tau$ for noise exposure on Delhi streets, using the sensor acquisition mode. A spatial sample of street segments is drawn from the road network, allocated into itineraries, and visited by a field team carrying a calibrated sound meter. Measurements are taken at every 5 km along the route and at every traffic stop. Each observation stores GPS, timestamp, and the minimum and maximum decibel readings within a fixed measurement window.

Across 55 logs collected over four fieldwork days and 12 itineraries, the median minimum reading is 65.9 dB and the median maximum 95.4 dB; the highest peak is 119.2 dB. The threshold-exceedance estimand $P_{55}$ for minimum readings is 87.3 percent: in most measurement windows, the lower envelope of street-level sound exceeds the daytime residential ambient standard under India's Noise Pollution Rules \citep{indiaNoiseRules2000}. The threshold-exceedance estimand $P_{100}$ for maximum readings is 30.9 percent. Traffic stops register about 8 dB higher at the floor of the measurement window and 5 dB higher at the peak than other sampled stops; treated as an opportunistic oversample of a high-exposure condition, they are reported separately rather than pooled into the citywide mean. The walk-based standard error around $\mu$ is the design-preferred standard error because the route is the primary sampling unit and within-route observations are systematic sub-samples.

The estimand structure differs from the imagery application. For noise, the relevant quantities are a population mean and a threshold-exceedance share rather than a per-unit ratio. The sampling and inference machinery is the same: probability sample on the road network, allocator-routed itineraries, design-based variance.

\subsection{Public goods and infrastructure}

The Street View pipeline annotates infrastructure features (footpaths, lane markings, potholes, litter, land use) on the same frames that produce the gender counts, and the sound meter pipeline records noise alongside contextual covariates such as road class, distance from the carriageway, and construction activity. Both support a regression of the form
\begin{equation}
    Y_i \;=\; \alpha + \beta_1 \text{Footpath}_i + \beta_2 \text{Litter}_i + \beta_3 \text{RoadType}_i + \gamma_c + \epsilon_i,
\end{equation}
where $Y_i$ is the per-unit female share or the per-unit decibel level, and $\gamma_c$ absorbs city or neighborhood fixed effects. The regression describes conditional associations between observable infrastructure and the local outcome. It is not causal in this design: infrastructure placement is not as-if random, and unobserved neighborhood characteristics may drive both the infrastructure and the outcome. The estimated coefficients identify which public-good deficits co-occur with low female presence and with high noise, a task that is informative for descriptive political analysis even when causal identification requires a separate research design.

\subsection{Targeting and program evaluation}

Design-based estimates produce defensible citywide quantities. A model trained on the sampled annotations or readings produces predicted-outcome maps for unsampled locations, which can guide operational targeting: where to send safety auditors, where to survey women about perceived safety, where to pilot lighting and sidewalk investments, where to deploy noise barriers or traffic-calming measures. The two outputs play different roles. The design-based estimate is the citywide claim. The model-based map is a prioritization tool whose predictions are conditional on the training sample and the predictor set.

The framework supports panel and difference-in-differences designs by repeating the sampling protocol over time. Sample the same segments before and after a policy change, or sample treatment and comparison areas with repeated time windows; the outcome is the change in $R$ or $\mu$. With randomized or plausibly exogenous rollout, the resulting estimator is causal under the standard parallel-trends or unconfoundedness assumptions. Without exogenous rollout, the same data support transparent event-study designs, provided that imagery dates, sensor calibration, and collection windows are explicitly controlled.

\section{Conclusion}

We develop a probability-sampling framework for street-level field measurement. The frame is constructed from the road network. Sampling is randomized over street segments or independent random walks. Allocation is operational, not inferential. Estimation produces design-based estimates of population ratios, means, and threshold-exceedance shares, with variance estimators for simple random, stratified, clustered, and walk-based designs. Coverage and nonresponse are part of the design.

Three acquisition modes specialize the framework: researcher-collected imagery, researcher-deployed sensors, and platform imagery. The first two preserve randomization over time at the cost of operational complexity. The third is scalable but constrained by face and plate obfuscation, by collection at platform-determined times rather than researcher-randomized times, and by image dates that can lag the question by years. Treating platform imagery as a measurement instrument with these specific properties, rather than as a substitute for researcher-collected data, makes the resulting biases auditable.

The applications show what the framework produces. Across Mumbai, Navi Mumbai, and Delhi, the people-weighted female share of visible pedestrians is roughly 18 to 24 percent, with predictable variation across road types and acquisition modes. In Delhi, the lower envelope of street-level noise exceeds the daytime residential ambient standard in 87 percent of measurement windows, and 31 percent of windows contain a peak at or above 100 dB. Both sets of estimates depend on a single design choice, that the sample is drawn from a known frame on the road network, with inclusion probabilities recorded and used.

\clearpage
\bibliographystyle{plainnat}
\bibliography{references}

\end{document}
